\chapter{实验结果与分析}\label{sec:experiments-cn}
\section{实验设置、模型归属与评价指标}\label{sec:experiment-settings-cn}
实验以产品回片完工时间$C_{\max}$为主指标，并分别记录有效下界、间隙、首次可行时间、求解器时间与工作进程墙钟。等待和节拍由解码后的物理事件重新计算。完整模型最优、受限子问题最优、限时有解和限时无解分别报告。

默认配置D采用100秒工艺腔非加工驻留与机械手承片上限，时域10000秒，LL复位尾部也在时域内。历史配置E采用10000秒模块总驻留与机械手承片上限，时域100000秒，未额外约束复位尾部；其机制结果单列。两者通常采用8片共享辅助晶圆、允许提前清洁及非递减出片规则；库存/清洁机制实验另列改变因素。两种合成工艺的单片累计处理时长均为180秒，各产品的模式与腔室资格由实例预先指定。

计算环境为Windows 11、20逻辑处理器、约32 GB内存，Python 3.12.12、PySCIPOpt 6.1.0，每次求解单线程。本章已完成的期刊补充实验不设固定SCIP/RSS配额，以系统可用内存低于10\%作为保护条件；其他批次按各自冻结资源设置说明。已有冻结结果保持原资源协议；模型、策略及输入解哈希随结果保存。三个求解种子或训练种子不作为三个独立实例。

正常分离的SPBS臂为60秒初始化改善加240秒主求解，其他臂为300秒主求解，未用完的60秒不再分配。该预算从共同的已验证输入开始计时，原始输入的外部构造成本单独记录。首轮根分离有一次300秒设置实际运行428秒，节拍搜索有一次600秒设置实际802秒，均不作为严格等时证据。跨训练种子补充组另有两次240秒主求解实际为245和243秒，按实际耗时保留；新六实例对照的414次主求解均未超过各自预算。
\section{模型编码与小规模最优性验证}
D的编码对照24次均通过检查。两个纯$4\times1$实例的两种表示均证明324秒最优；同LP实例变量数由3225降至3219，平均求解时间由55.33降至21.00秒。2片混合实例中Pjob与逐片表示的平均可行上界分别为518.00和514.33秒，4片混合均为542秒；因此表示压缩与搜索表现具有不同的变化趋势。以下8个小规模实例完整列出4次最优、1次限时可行及3次无可行解；无可行解不表示不可行。
\begin{table}[!htbp]
\centering\small
\caption{默认配置D的小规模实例求解结果（单位：s）}
\label{tab:default-small-new-cn}

\begin{adjustbox}{max width=\linewidth}
\begin{tabular}{llll}\toprule
$N$/配方/$H$ & LP分布 & 下界/上界 & 等待：前$\to$后\\\midrule
1/$4\times1$/3 & LP1 & 296.00/296.00 & $198\to6$\\
1/$2\times2$/2 & LP1 & 303.00/303.00 & $642\to54$\\
2/$4\times1$/3 & 跨LP & 324.00/324.00 & $72\to47$\\
2/$4\times1$/3 & 同LP & 324.00/324.00 & $244\to56$\\
2/混合/3 & 跨LP & 296.00/518.00 & $207\to55$\\
2/$2\times2$/2 & 跨LP & 207.10/-- & --\\
2/$2\times2$/2 & 同LP & 293.50/-- & --\\
4/混合/3 & 跨LP & 116.00/-- & --\\
\bottomrule\end{tabular}
\end{adjustbox}
\par\smallskip\footnotesize 前四行主问题最优，第五行可行未证最优，后三行超时无可行解。``--''不表示不可行；精修主目标容差为$10^{-6}$~s。
\end{table}
\section{结构化初始化与HEM组件消融}\label{sec:journal-ablation-cn}\label{sec:algorithm-contribution-cn}
\subsection{跨训练种子的初始化对照}\label{sec:journal-training-cn}
四个已有完整可行输入实例分别为2片纯$4\times1$跨LP/同LP、2片混合/$H=3$与4片混合/$H=10$。正常分离初轮72次和补充96次合计168次，全部赋值通过完整D模型检查及物理重放。HEM和RL-DMCT-S使用三个冻结训练种子、三个求解种子；SCIP只有三个求解种子，不人为复制训练重复。四个实例内各方法共享相同原始输入，初始完工依次为324、324、518和542秒。
\begin{table}[!htbp]
\centering\small
\caption{正常分离下跨训练种子的初始化对照（产品完工均值，s）}
\label{tab:journal-training-cn}
\begin{adjustbox}{max width=\linewidth}
\begin{tabular}{llrr}\toprule
实例 & 方法 & 无SPBS & 有SPBS\\\midrule
2片混合 & SCIP & 516.67 & 509.00\\
 & HEM & 517.14 & 508.89\\
 & RL-DMCT-S & 517.44 & 507.67\\
4片混合 & SCIP & 542.00 & 541.00\\
 & HEM & 542.00 & 541.00\\
 & RL-DMCT-S & 542.00 & 541.00\\
\bottomrule\end{tabular}
\end{adjustbox}
\par\smallskip\footnotesize 学习方法每格为三个训练种子与三个求解种子的9次均值，SCIP为3次均值。其余两个纯$4\times1$实例所有臂均保持已证最优的324~s。所有方法均从已有可行输入出发。
\end{table}
2片混合中，SPBS的三份起点为511、505、511秒，各训练种子相同。RL-DMCT的训练种子均值为505、509、509秒，HEM+SPBS为508.67、509、509秒，学习阶段的额外收益集中在训练种子1。4片混合均从542降至541秒，收益来自SPBS。四实例上RL-DMCT相对HEM+SPBS的平均相对改善约0.0600\%，胜/平/负为1/3/0，增量集中在2片混合实例和训练种子1。
\subsection{选割组件与分离协议}
根分离的九臂覆盖108次，训练种子固定为1，三个求解种子。前三种方法分别比较有/无SPBS，另加入特征扩展、去重排和贪心替代束搜索。正常分离六臂首轮72次包含在前述168次中。下表列出无SPBS组件对照。
\begin{table}[!htbp]
\centering\small
\caption{根分离协议的无SPBS选割组件对照（产品完工均值，s）}
\label{tab:journal-components-cn}

\begin{adjustbox}{max width=\linewidth}
\begin{tabular}{lrrrr}\toprule
方法 & 跨LP纯 & 同LP纯 & 2片混合 & 4片混合\\\midrule
HEM & 324 & 324 & 514.02 & 542\\
仅扩展特征 & 324 & 324 & 515.33 & 542\\
去结构重排 & 324 & 324 & 516.00 & 542\\
贪心替代束搜索 & 324 & 324 & 511.00 & 542\\
RL-DMCT-S & 324 & 324 & 512.00 & 542\\
\bottomrule\end{tabular}
\end{adjustbox}
\end{table}
在2片混合实例中，RL-DMCT-S的均值512.00秒低于去重排臂的516.00秒，高于贪心解码臂的511.00秒；特征扩展臂515.33秒高于HEM的514.02秒。正常分离增加学习回调后，无SPBS臂的平均上界仍高于根分离结果。加SPBS时，根分离中HEM的均值较低，正常分离中RL-DMCT的均值较低。初始化的改善在两种协议下均出现，学习组件的增量则取决于所用检查点和分离设置。
\subsection{固定结构与后续联合搜索的互补作用}
同一四实例、三求解种子比较联合300秒、固定结构300秒、先固定60秒再联合240秒，共36次，均经完整D验证。
\begin{table}[!htbp]
\centering\small
\caption{默认配置D的结构对照（三种子平均产品完工，s）}
\label{tab:journal-structure-cn}
\begin{adjustbox}{max width=\linewidth}
\begin{tabular}{lrrr}\toprule
实例 & 联合 & 固定 & 先固定后联合\\\midrule
2片/$4\times1$/跨LP & 324 & 324 & 324\\
2片/$4\times1$/同LP & 324 & 324 & 324\\
2片/混合/$H=3$ & 513.34 & 507.00 & 505.67\\
4片/混合/$H=10$ & 542 & 542 & 542\\
\bottomrule\end{tabular}
\end{adjustbox}
\end{table}
2片混合实例中，固定结构得到507秒的受限最优值，开放结构后进一步获得503秒可行解，显示扩大离散选择范围可在已有起点上取得额外改善。其余三实例目标无差异。
\subsection{扩展驻留配置下的结构搜索对照}
\label{sec:structure-comparison-cn}

该组采用历史扩展驻留配置E，检验有限预算内利用已有排程结构降低产品完工上界的作用。其结果与默认配置D的结构对照分别报告。

选择$N\in\{4,8\}$、$H\in\{3,10\}$的四个双源混合实例，比较联合搜索、固定初始结构以及先固定后联合三种策略，共36次运行。每个实例和种子的三种策略采用同一初始排程，
总主问题求解预算均为300~s；先固定后联合分配60~s与240~s。离线初始化与模型构建时间单独记录。表~\ref{tab:structure-ablation-cn}以产品完工时间作为比较指标。

\begin{table}[!htbp]
\centering\small
\caption{四个混合实例的搜索策略对照（三种子均值，单位：s）}
\label{tab:structure-ablation-cn}
\begin{tabular}{ccccc}\toprule
$N$ & $H$ & 联合搜索 & 固定结构 & 先固定后联合\\\midrule
4 & 3 & 854.53 & 543.67 & 647.00\\
4 & 10 & 854.53 & 542.00 & 542.00\\
8 & 3 & 2071.00 & 2029.33 & 2071.00\\
8 & 10 & 1720.00 & 1249.33 & 1575.00\\
\bottomrule\end{tabular}
\end{table}

如表~\ref{tab:structure-ablation-cn}所示，固定结构在四个实例中均取得低于直接联合搜索的平均可行上界，其中8片、$H=10$实例由1720.00降至1249.33~s。先固定后联合仅在4片、$H=10$实例上与固定结构持平，其余三例的平均上界更高。这表明，在本组300~s预算内，围绕已有结构优化时序比扩大离散搜索范围更有效；联合模型同时提供完整问题的有效下界。


\section{清洁、加工模式与有限辅助库存}\label{sec:journal-mechanism-cn}
24配置来自六种工作负载、两个清洁阈值$H\in\{2,3\}$和两档库存$D\in\{4,8\}$；各比较提前清洁与原触发规则、三个求解种子，共144次。两规则使用同一完整候选域，共用相同有效输入或共同冷启动，主求解600秒且关闭精修。全量结果为44最优、9限时可行和91限时无可行解，53份有效排程均通过完整模型与物理重放。八个单产品配置在两规则下均存在最优证书。
\begin{table}[!htbp]
\centering\small
\caption{单产品场景的最优产品完工时间（s）}
\label{tab:journal-mechanism-cn}
\begin{adjustbox}{max width=\linewidth}
\begin{tabular}{rrrr}\toprule
$H$ & 辅助片数 & 纯$4\times1$ & 纯$2\times2$\\\midrule
2 & 4 & 296 & 309\\
2 & 8 & 296 & 303\\
3 & 4 & 296 & 309\\
3 & 8 & 296 & 303\\
\bottomrule\end{tabular}
\end{adjustbox}
\par\smallskip\footnotesize 每个表值在允许提前清洁与原触发规则下相同，并由相应配置的有效上下界确认；这是配置最优值，不是各次限时运行的均值。两种合成工艺的单片累计处理时长均为180~s。
\end{table}
单产品$2\times2$场景的辅助库存由4片增至8片，使最优完工缩短6秒（1.94\%），而$4\times1$保持296秒。$H=2$的链式场景有11份连续加工、满片清洁及身份复用同时存在的有效排程，其中4片库存的五份排程以四个身份完成9次借用，即5次重复借用。阈值由2改成3使单产品场景的清洁次数由1变为0，但相应库存的产品完工最优值不变。

2片混合、$H=3$、8片辅助库存场景中，提前清洁与原规则的最好可行值分别为511和514秒，对应最大有效下界为233.93和144.00秒，均保留较大界差。其余多产品配置均在时限内未取得可行解，困难主要出现在初始排程的获取阶段。
\subsection{连续加工、满片清洁与重复借用见证}\label{sec:d-coupled-witness-cn}
既有D下单产品$2\times2$/$H=2$、8片库存的排程达到完整主问题最优303秒。产品所在PW对连续参与两个周期，随后四片辅助片满片清洁。9次借用使用7个实体身份；固定结构精修使总等待由642降至54秒，全部占用结束由978提前至796秒，产品完工保持303秒。该例展示连续加工、满片清洁与重复借用的联合可行排程；库存的性能影响由前述4片与8片配对实验给出。
\begin{figure}[!htbp]
\centering
\includegraphics[width=\linewidth]{fig_coupled_schedule_merged_cn}
\caption{产品晶圆连续加工及辅助片借用区间。上行为产品晶圆的两次$2\times2$加工；D1--D7的蓝色、绿色区间分别表示生产、清洁用途的借用。虚线和点线分别表示303~s产品完工与796~s全部占用结束。}
\label{fig:d-coupled-witness-cn}
\end{figure}
\section{两阶段时序精修与节拍消融（配置E）}\label{sec:stage2-validation-cn}
配置E的30组阶段内配对均通过模型检查和物理重放；第一阶段270秒、精修30秒，同一次运行的前后解比较，全部保持产品完工上界。普通资格24次总等待降幅29.66\%至95.85\%，受限资格6次46.74\%至67.40\%；该指标累加产品和辅助晶圆的受计量路线等待，反映固定产品完工时间下的时序压缩。
\begin{table}[!htbp]
\centering\small
\caption{配置E同一排程的时序精修（三种子均值，单位：s）}
\label{tab:refinement-ablation-cn}
\resizebox{\columnwidth}{!}{%
\begin{tabular}{ccccc}\toprule
$N/H$ & 供片来源与腔室限制 & $C_{\max}$ & 精修前等待 & 精修后等待\\\midrule
4/3 & 单源 & 1491.00 & 1648.00 & 840.00\\
4/3 & 双源 & 854.53 & 3991.73 & 333.67\\
4/10 & 单源 & 1491.00 & 1185.00 & 600.00\\
4/10 & 双源 & 851.97 & 1259.65 & 626.67\\
8/3 & 单源 & 3933.00 & 3158.00 & 1680.00\\
8/3 & 双源 & 2071.00 & 10943.00 & 536.00\\
8/10 & 单源 & 2985.00 & 2276.00 & 1200.00\\
8/10 & 双源 & 1720.00 & 10998.02 & 692.00\\
8/3 & 单源、限制腔室 & 4707.00 & 3830.00 & 2040.00\\
8/3 & 双源、限制腔室 & 2264.00 & 1046.00 & 341.00\\
\bottomrule\end{tabular}}
\end{table}
\begin{table}[!htbp]
\centering\small
\caption{节拍权重消融（配置E，三种子均值，时间单位：s）}
\label{tab:rhythm-weight-cn}
\begin{adjustbox}{max width=\linewidth}
\begin{tabular}{lrrrr}\toprule
实例 & $\lambda$ & $C_{\max}$ & 总等待 & 节拍偏差\\\midrule
16片纯$2\times2$ & 0 & 2537 & 1000 & 792\\
 & 5 & 2537 & 1016 & 784\\
 & 25 & 2537 & 1016 & 784\\
 & 100 & 2537 & 1016 & 784\\
24片混合 & 0 & 5199 & 1440 & 20\\
 & 5 & 5199 & 1440 & 20\\
 & 25 & 5199 & 1440 & 20\\
 & 100 & 5199 & 1440 & 20\\
\bottomrule\end{tabular}
\end{adjustbox}
\par\smallskip\footnotesize 两实例分别含6个和2个有效节拍中心；各臂最大等待均为65~s。
\end{table}
上表的16/24片两个实例分别有6/2个节拍中心，三种子共享各自第一阶段结构，只改变权重；每臂60秒，24次均达到受限精修最优。16片实例的偏差由792降至784秒，总等待由1000增至1016秒，权重5、25和100的指标相同；24片实例在四档权重下均无变化。这一结果体现了等待压缩与节拍均衡的折中。D下另外9次要求至少两个节拍中心的搜索均未取得可行输入，后续权重消融未执行。
\section{供片与库存变化下的初始化和选割收益}
\label{sec:journal-transfer-cn}

\begin{table}[!t]
\centering\small
\caption{不同配方与供片场景下的算法对照（产品完工均值，s）}
\label{tab:journal-transfer-cn}
\begin{adjustbox}{max width=\linewidth}
\begin{tabular}{cccccccc}\toprule
配方 & 场景 & SCIP & HEM & RL-DMCT-S & SCIP+SPBS & HEM+SPBS & RL-DMCT\\\midrule
$4\times1$ & S1 & 324.00 & 324.00 & 324.00 & 324.00 & 324.00 & 324.00\\
$4\times1$ & S2 & 324.00 & 324.00 & 324.00 & 324.00 & 324.00 & 324.00\\
$4\times1$ & S3 & 324.00 & 324.00 & 324.00 & 324.00 & 324.00 & 324.00\\
$2\times2$ & S2 & 342.00 & 342.67 & 342.00 & 344.00 & 344.00 & 343.33\\
混合 & S2 & 521.00 & 520.00 & 520.11 & 492.00 & 492.00 & 492.00\\
混合 & S3 & 327.00 & 327.00 & 326.67 & 327.00 & 327.00 & 326.11\\
\bottomrule\end{tabular}
\end{adjustbox}
\end{table}

为检验初始化和选割在不同输入条件下的效果，本节改变供片分布与辅助库存，分别比较初始排程的改善，以及各选割组件从相同初始排程出发获得的进一步改善。

\paragraph{供片、清洁阈值与库存变化下的共同起点}

为考察不同配方与供片场景下的算法表现，将$N\in\{2,4,8\}$、三种配方与三个场景交叉，形成27个测试实例。场景S1将产品集中于LP1、$H=2$、辅助库存为4片；S2按LP1--LP4轮转分配产品、$H=3$、辅助库存为6片；S3将前半产品置于LP1、后半置于LP2、$H=5$、辅助库存为8片。各场景通过改变供片分布、清洁频率和辅助库存覆盖不同输入组合。使用SCIP在无初始解的条件下搜索可行排程，求解种子为17、名义预算为1200~s，得到六个2片实例的可行排程，资格覆盖为6/27，其余21例在预算内未取得可行解。后续算法对照使用这六份排程作为共同起点；初始解搜索时间与后续300~s改善预算分别统计。

在这六个2片实例上，采用正常分离，沿用多训练种子对照并增加三种带SPBS的消融变体，共414次算法运行，加上18次共享准备形成432个队列单元；432份赋值均通过完整模型约束检查和独立物理重放。每个实例与求解种子组合执行一次共享SPBS，共得到18份初始排程，其中3份降低产品完工时间；同组所有带SPBS的方法使用相同排程。209次运行证明完整主问题最优，205次取得限时可行解，实际主求解均未超过设定预算。

\paragraph{初始化收益随输入结构和预算分配变化}

如表~\ref{tab:journal-transfer-cn}所示，三个纯$4\times1$实例在各方法下均证明324~s最优。混合S2的原输入为521~s，共享SPBS分别给出494、494、488~s，均值492~s；各方法均保留这些起点，因此该实例的改善来自初始化。混合S3的共享起点均为327~s，RL-DMCT进一步将均值降至326.11~s，而HEM加SPBS保持327~s。纯$2\times2$ S2上，RL-DMCT的343.33~s低于HEM加SPBS的344~s，但高于无SPBS时的342~s；该实例初始化未改善起点，同时缩短了后续主求解时间。结合第~\ref{sec:journal-training-cn}节的结果，SPBS对部分混合实例形成了有效起点，其收益随输入结构及预算分配而变化。

\paragraph{共享初始解后选割组件带来的额外改善}

表~\ref{tab:journal-transfer-components-cn}给出了共享SPBS输入后的组件对照。对RL-DMCT与HEM加SPBS按同一实例、训练种子和求解种子配对，先在实例内平均相对完工时间差，再对六个实例等权汇总，平均改善为0.0776\%，中位数为0，胜/平/负为2/4/0；按实例重采样的描述性95\%区间为$[0,0.1682]\%$。三个训练种子的实例平均改善依次为0.0510\%、0.1309\%和0.0510\%。

\begin{table}[!htbp]
\centering\small
\caption{共享SPBS输入的组件对照（产品完工均值，s）}
\label{tab:journal-transfer-components-cn}
\resizebox{\columnwidth}{!}{%
\begin{tabular}{ccccc}\toprule
方法 & 三个纯$4\times1$ & $2\times2$ S2 & 混合S2 & 混合S3\\\midrule
HEM+SPBS & 324.00 & 344.00 & 492.00 & 327.00\\
仅扩展特征+SPBS & 324.00 & 344.00 & 492.00 & 326.67\\
去结构重排+SPBS & 324.00 & 344.00 & 492.00 & 327.00\\
贪心替代束搜索+SPBS & 324.00 & 344.00 & 492.00 & 326.67\\
RL-DMCT & 324.00 & 343.33 & 492.00 & 326.11\\
\bottomrule\end{tabular}}
\end{table}

表~\ref{tab:journal-transfer-components-cn}中，完整组合在纯$2\times2$ S2和混合S3上的平均完工时间均低于三种消融变体，其余四实例持平。相较仅扩展特征、去结构重排和贪心替代束搜索，六实例平均配对改善分别为0.0605\%、0.0776\%和0.0606\%。结合第~\ref{sec:journal-ablation-cn}节根分离下贪心替代束搜索结果更好的对照可见，完整选割组合在共享起点后取得了额外改善，束搜索的收益则依赖求解设置。


\section{满片加工与共享资源约束的排程验证}
\label{sec:coverage-scale-cn}

本节通过两份完整排程展示设备执行过程：2片实例的最优排程说明辅助片如何补足满片加工所需工位，16片实例的可行排程展示两腔如何共用搬运资源。两份排程均通过完整模型约束检查和独立事件时序验证。

\subsection{2片最优排程中的辅助片补位}

图~\ref{fig:small-optimal-gantt-cn}给出2片同LP纯$4\times1$实例的一份产品完工时间最优排程。两片产品均属于Pjob1且由LP1供片，辅助库存为8片，清洁阈值为3。主问题下界与排程的产品完工时间均为324~s，因此该排程在产品完工时间目标下达到全局最优。

\begin{figure}[!htbp]
\centering
\includegraphics[width=\columnwidth]{fig_schedule_n2_cn}
\caption{2片同LP纯$4\times1$产品的完工时间最优排程。P1[1:2]表示Pjob1的两片产品，D1、D2为辅助片。虚线与点线分别标示324~s产品回片和342~s全部占用结束。}
\label{fig:small-optimal-gantt-cn}
\end{figure}

该排程以两片产品组成一个PW对，以辅助片D1、D2组成另一个PW对，共同填满CH2的四个工位，并在104--284~s完成180~s加工。产品片依次经过ATR搬运、对准器和上装载锁，由VTR装入腔室；辅助PW对先于产品PW对装入，使加工开始时满足满片条件。加工结束后，产品经下装载锁返回LP1，辅助片于294~s归还真空辅助库。

两片产品均于324~s回到LP1，下装载锁的压力复位持续至342~s。该排程展示了产品数量不足四片时的辅助片补位与完整取还流程，并区分了产品回片和设备复位两个结束时刻。腔室仅执行一个加工周期，未达到清洁阈值。

\subsection{16片可行排程中的双腔加工与搬运安排}

图~\ref{fig:current-defaults-gantt-cn}展示16片纯$4\times1$产品的可行排程。四个LP各含4片产品，均属于Pjob1并可进入CH2与CH3；辅助库存为8片，清洁阈值为3。

\begin{figure}[!htbp]
\centering
\includegraphics[width=\columnwidth]{fig_schedule_n16_cn}
\caption{16片纯$4\times1$产品的完整资源可行排程。LP1--LP4各含4片产品，均属于Pjob1。虚线与点线分别标示920~s产品回片和938~s全部占用结束。}
\label{fig:current-defaults-gantt-cn}
\end{figure}

该排程以预先构造的晶圆动作序列为输入，将其配对、腔室分配及事件时刻代入完整模型进行可行性检验。模型包含64个候选PW对、每腔32个候选周期，共177905个变量与671864条约束。全部赋值满足产品可加工腔室限制、有限资源容量、非加工驻留及同任务非递减投片约束，独立事件时序验证也通过。

该排程的产品完工时间为$C_{\max}=920$~s，16片产品组成8个PW对，共执行4个满片周期。CH2于170--350~s与497--677~s分别加工Pjob1的第1--4片和第9--12片；CH3于326--506~s与645--825~s分别加工第5--8片和第13--16片。每片均完成一次180~s工艺处理，按已验证的搬运与装载锁流程返回原LP。

图~\ref{fig:current-defaults-gantt-cn}给出152个事件的时间安排，包括LP取还、ATR、对准器、上下装载锁、VTR以及两工艺腔的占用。腔室加工可与其他产品的供片和回片搬运同时进行，各PW对的搬运次序受共享机械手与装载锁约束。全部产品于920~s回到LP，下装载锁的压力复位持续至938~s，产品完工与全部占用结束相差18~s。每腔仅加工两个周期，尚未达到清洁阈值3；全部周期由产品片满载，未借用辅助片。该排程展示了满片条件下两腔加工与共享资源的使用次序。


\FloatBarrier
\section{100片完整排程与模型适用边界}
历史固定结构模型保存了100片混合配方的有效排程，其中50片为$4\times1$、50片为$2\times2$，产品完工时间为9998秒。该结果保存了原模型独立赋值检查，属于可行解，未证明全局最优。其工艺计时、辅助片配置与候选结构不同于本章默认配置D，因此仅作为既有大规模排程的历史示例，不纳入D的性能排名。默认配置D的75片混合批次见第~\ref{sec:pareto-evidence}~节，100片完整对照仍待验证。

图~\ref{fig:legacy100a}和图~\ref{fig:legacy100b}按产品编号分两页展示全部100片的轨迹，附图保留原始SVG中各动作的起止时刻；完整辅助片与资源视角随论文资料包保存。
\begin{figure}[p]\centering\includegraphics[width=\linewidth,height=0.84\textheight,keepaspectratio]{legacy100_products_1}\caption{历史固定结构模型100片混合排程：产品1至50（全部时间轴）}\label{fig:legacy100a}\end{figure}
\begin{figure}[p]\centering\includegraphics[width=\linewidth,height=0.84\textheight,keepaspectratio]{legacy100_products_2}\caption{历史固定结构模型100片混合排程：产品51至100（全部时间轴）}\label{fig:legacy100b}\end{figure}
\FloatBarrier

\clearpage
\section{双目标前沿、组件消融与75片排程}\label{sec:pareto-evidence}

\subsection{两目标的冲突与可认证前沿}

两片算例S1含同一Pjob的两片$2\times2$产品、四片辅助片，清洁阈值为2。采用该工况的配对与双腔时序子模型，在82.03~s内重新生成图~\ref{fig:pareto-s1}的七份排程。逐点重放并与独立数值证书比对，在$2\times10^{-5}$容差内命中七个全局非支配点，含三个非支持点。最快回片为331~s，对应$B=424/3$。最低$B=42$对应519~s。当上限为$B\le100$时，两片分别与辅助片配对，最短工期为514~s。只从线性加权和的支持点中选解则需519~s。该例的$R=0$，两目标的冲突来自交付与路线等待。专用重算的HV为0.378339，$\mathrm{IGD}^{+}=2.45\times10^{-10}$。通用领域搜索在40~s内每个种子命中一个认证点，种子1延长至300~s后命中两个点。
\begin{figure}[!htbp]
\centering\includegraphics[width=.92\linewidth]{pareto_s1}
\caption{两片算例S1在两种配对方式下的完整前沿。两图分别标注原始坐标，三枚菱形为线性加权和无法取得的非支持点。}\label{fig:pareto-s1}
\end{figure}

四片标称混合工况M4中，发现了413~s交付的排程。已证最低$B$为54，达到该值至少需要476~s，已验证排程将对应的最短工期限定在$[476,497]$~s。因此追求最低$B$至少要比413~s的排程多等63~s。在已访问的35种整数结构中，开放连续时刻后有一段位于这些结构合并边界上的有效线段，其工期区间为$(420.5,423)$~s。该区间的结论限于这35种结构。

\subsection{等预算对照与扰动}

五个工况各比较六种方法，每个方法使用三个种子。S1、M4、H4、H8和H16的产品组成（$4\times1$片数，$2\times2$片数）依次为$(0,2)$、$(2,2)$、$(3,1)$、$(4,4)$和$(8,8)$。同工况、同种子的各臂共享可行起点。NSGA-II采用非支配排序与拥挤距离选择\citep{deb2002}，MOEA/D采用切比雪夫分解与邻近权重子问题替换\citep{zhangli2007}。两种基线以Python实现并共用设备解码器。$\varepsilon$-SCIP直接在原模型中扫描目标上限，原两阶段法保持工期优先。表~\ref{tab:pareto-small}给出各方法返回集合的超体积（HV）中位数。对归一化的双目标计算该集合支配至参考点$(1.1,1.1)$的面积，数值越大表示覆盖越广。S1按认证前沿极值定标，其他工况按共同起点$(A_0,B_0)$取$[0,1.5A_0]\times[0,2B_0]$。HV在同一工况内比较。
\begin{table}[!htbp]\centering\small
\caption{五个工况的三种子HV中位数与单次搜索时限}\label{tab:pareto-small}
\begin{tabular}{lrrrrrrr}\toprule
工况 & 时限/s & 领域搜索 & 加选割 & $\varepsilon$-SCIP & NSGA-II & MOEA/D & 两阶段 \\\midrule
S1 & 40 & 0.27468 & 0.27468 & 0.00000 & 0.27468 & 0.27468 & 0.00000 \\
M4 & 45 & 0.54479 & 0.54116 & 0.26126 & 0.53952 & 0.53947 & 0.35375 \\
H4 & 45 & 0.75582 & 0.75448 & 0.26000 & 0.70736 & 0.70690 & 0.35386 \\
H8 & 55 & 0.58856 & 0.58856 & 0.26000 & 0.57961 & 0.61188 & 0.28050 \\
H16 & 80 & 0.37344 & 0.38830 & 0.26000 & 0.47203 & 0.50191 & 0.26000 \\
\bottomrule\end{tabular}
\end{table}

领域搜索在两种四片混合工况的HV中位数均最高。八片、十六片工况由MOEA/D取得最高中位数。用同一随机种子配对比较，M4开启配对下界、在已有下界时采用领域邻域、允许改变对内取放顺序，对应的HV差值中位数分别为$+0.00372$、$+0.01710$和$+0.06650$，三项改动在三个种子中均提高HV。自适应选择上限的差值为$-0.02102$，接入冻结选割策略的差值为$-0.00330$。在12对固定上限、共享起点的4~s原模型子问题中，下界对其中6对排除了至少一种配对数量，12对的条件工期均相同。

M4另将邻域释放比例从0.30分别改为0.15和0.50，两个配对种子的HV差值中位数分别为$-0.03107$和$-0.02767$。把修复预算比例从0.22改为0.12和0.32，对应变化为$-0.00734$和$+0.00105$。在八片节拍实验R8中，151份归档排程有46份产生非零节拍偏差。对同批候选在相同工期上限下分别按0、12.5、25和50重算节拍权重，五个有可选排程的种子与工期上限组合均选择了相同的$R$值。这些数据区分了节拍项实际激活与改变权重后决策发生变化两种情况。

\subsection{75片混合批次的可行前沿}

采用默认配置D，37片执行$4\times1$，38片执行$2\times2$。LP1至LP4的两类产品数量分别为$(10,9)$、$(9,10)$、$(9,10)$和$(9,9)$。两个Pjob共用CH2和CH3，真空辅助库存为8片，清洁阈值为3，时域为10000~s。历史构造得到的可行起点为$(A_0,B_0)=(8630,3932)$。四种方法由同一起点出发，各运行三个种子，每次限时90~s。共享起点的外部准备与完整原模型的事后复核单独计时，搜索时限包含结构调整、重定时和设备事件复验。固定结构对照仅改变动作时刻及$B$上限，在预设网格的扫描区间耗尽后提前停止。

每次运行按预设规则提交至多七点，先保留两个端点，再依次选择与已保留集合的归一化距离最大的点。表~\ref{tab:pareto75}的HV仅使用当次提交点。最后两列分别对各次提交集合的最小$A$与最小$B$取中位数，两列可以对应不同的排程。
\begin{table}[!htbp]\centering\small
\caption{75片混合批次的三种子结果。HV范围为三次运行的最小值至最大值。}\label{tab:pareto75}
\begin{tabular}{lrrrr}\toprule
方法 & HV中位数 & HV范围 & 最小$A$/s & 最小$B$/加权秒 \\\midrule
领域搜索 & 0.44494 & 0.44391--0.44546 & 7623 & 1786 \\
NSGA-II & 0.43287 & 0.41437--0.43561 & 7586.88 & 1951 \\
MOEA/D & 0.43898 & 0.43789--0.44693 & 7662 & 1851.5 \\
固定结构扫描 & 0.27137 & 0.27137--0.27137 & 8630 & 3722 \\
\bottomrule\end{tabular}
\end{table}

领域搜索的HV中位数为0.44494，MOEA/D为0.43898。第二个种子中MOEA/D取得0.44693，高于同种子的领域搜索。固定生产块和资源次序后，单靠重定时仍能在$(8630,3932)$与$(8852,3722)$之间找到取舍。开放结构后的三个搜索方法均改善了固定结构法的中位HV。75片领域搜索在本组所用上限上未触发配对下界筛除，表中差异不能归因于该下界或学习选割。

图~\ref{fig:pareto75}汇集全部12次运行在限时内发现的非支配解，含单次提交时被七点上限省略的内部点。按两目标再次去支配后共得20个实际坐标。在最短工期端，$(A,B)=(7567,2599)$。在最低等待与节拍指标端，$(A,B)=(8045,1782)$。从前者换到后者，批次晚交付478~s，$B$减少817加权秒。若只在本次候选中选择，要求$B\le2000$时最早可在7639~s交付。该值可与两端点一起用于选择交期。
\begin{table}[!htbp]\centering\small
\caption{75片共同起点与按交付要求选取的三份排程。$A$、$L_{\rm sum}$、$L_{\max}$和$R$的单位为秒，$B$为加权秒。}\label{tab:pareto75-policies}
\begin{tabular}{lrrrrr}\toprule
排程选择 & $A$ & $B$ & $L_{\rm sum}$ & $L_{\max}$ & $R$ \\\midrule
共同起点 & 8630 & 3932 & 3312 & 62 & 0 \\
发现的最短工期 & 7567 & 2599 & 1884 & 49 & 9 \\
$B\le2000$时最早交付 & 7639 & 1951 & 1776 & 17.5 & 0 \\
发现的最低$B$ & 8045 & 1782 & 1612 & 17 & 0 \\\bottomrule
\end{tabular}
\end{table}
\FloatBarrier
\begin{figure}[!htbp]
\centering\includegraphics[width=.96\linewidth]{pareto_75}
\caption{75片混合批次的经验前沿。左图叠加各方法提交点、共同起点与全部发现记录的合并非支配点，右图放大最短工期端。每个非支配点均有独立核验的完整可行排程。}\label{fig:pareto75}
\end{figure}

两端点的排程均含九个可计量的节拍位置。沿最短工期端到最低$B$端，总等待由1884降至1612~s，最大单项等待由49降至17~s，节拍偏差由9降至0~s。三项分别解释$B$下降量中的272、320和225加权秒。

75片的12次运行共提交78份排程。对提交集合及补全后的合并前沿去重，共有63份排程逐片重放，并向原模型的6\,764\,011个变量赋值、逐一检查24\,271\,470条约束。按事件复算的$A$、$B$与原模型赋值相符。此前小规模主比较连同消融、参数扰动和长预算实验共164组，699份最终排程均已核对。小规模的完整性证书仅适用于S1。75片的网格、有限结构邻域与搜索预算限定了本图的覆盖范围。对完整可行域尚无排除其他非支配点的证书。
